\chapter{Resultats}

Un cop desenvolupat i executat el programa, vaig obtenir una sèrie de resultats que mostren l’estat de les estacions del Bicing en el moment de la consulta. En aquest apartat presento exemples concrets, taules amb dades i una anàlisi general del que es pot observar.

\section{Exemple de sortida del programa}
El primer resultat que s’obté és una llista de totes les estacions amb la seva informació bàsica: bicicletes lliures, espais lliures i canvis ràpids possibles. A continuació es mostra un fragment real de l’execució:

\begin{verbatim}
Estació: Gran Via - Marina
   -> Bicis lliures: 5
   -> Espais lliures: 7
   -> Canvis ràpids possibles: 5
----------------------------------------
Estació: Passeig de Gràcia - Casp
   -> Bicis lliures: 0
   -> Espais lliures: 12
   -> Canvis ràpids possibles: 0
----------------------------------------
Estació: Universitat - Pelai
   -> Bicis lliures: 8
   -> Espais lliures: 3
   -> Canvis ràpids possibles: 3
----------------------------------------
\end{verbatim}

Aquest fragment ja ens permet veure tres situacions diferents:
\begin{itemize}
    \item Estacions equilibrades (Gran Via - Marina), on hi ha prou bicis i espais.
    \item Estacions saturades d’espais buits però sense bicicletes (Passeig de Gràcia - Casp).
    \item Estacions amb moltes bicicletes però pocs espais (Universitat - Pelai).
\end{itemize}

\section{Taula de resultats}
Per fer més visual la informació, vaig resumir els resultats d’algunes estacions en una taula:

\begin{center}
\begin{tabular}{|l|c|c|c|}
\hline
\textbf{Estació} & \textbf{Bicis lliures} & \textbf{Espais lliures} & \textbf{Canvis ràpids} \\
\hline
Gran Via - Marina & 5 & 7 & 5 \\
Passeig de Gràcia - Casp & 0 & 12 & 0 \\
Universitat - Pelai & 8 & 3 & 3 \\
Plaça Espanya - Creu Coberta & 2 & 9 & 2 \\
Arc de Triomf - Ribes & 6 & 6 & 6 \\
\hline
\end{tabular}
\end{center}

Aquesta taula mostra clarament com el nombre de canvis ràpids depèn del recurs més limitat: a la Universitat - Pelai són només 3 malgrat tenir 8 bicicletes disponibles, ja que només hi ha 3 espais lliures.

\section{Anàlisi dels resultats}
A partir de les dades obtingudes, es poden fer diverses observacions:
\begin{itemize}
    \item Hi ha estacions molt equilibrades, amb un nombre semblant de bicis i espais, que permeten un gran nombre de canvis ràpids.
    \item Algunes estacions del centre tendeixen a tenir pocs espais lliures perquè hi ha molta demanda d’aparcar bicicletes.
    \item Altres estacions, especialment en zones més allunyades, presenten l’efecte contrari: molts espais lliures i poques bicicletes.
    \item El concepte de canvi ràpid resulta útil perquè resumeix la situació de cada estació en un sol número.
\end{itemize}

\section{Limitacions dels resultats}
Cal tenir en compte que aquests resultats són d’un moment concret i poden variar molt segons l’hora del dia o el dia de la setmana. Per obtenir una anàlisi més profunda, caldria registrar dades durant més temps i observar-ne l’evolució.

Tot i això, els exemples presentats demostren que el programa funciona correctament i és capaç de calcular els canvis ràpids a partir de les dades reals de l’API.
